% !TeX program = xelatex
\documentclass[12pt,a4paper]{article}
\usepackage[margin=2cm]{geometry}
\usepackage{fontspec}
\usepackage{amsmath,amssymb}
\usepackage{enumitem}
\setmainfont{Times New Roman}
\setlength{\parindent}{0pt}
\setlength{\parskip}{6pt}
\begin{document}
\begin{center}
{\Large\bfseries ĐÁP ÁN ĐỀ THI TUYỂN SINH VÀO LỚP 10 THPT}\\
{\bfseries Thành phố Hà Nội, năm học 2009--2010}\\
{\itshape Phần Đại số}
\end{center}

\textbf{Câu 1.}
\begin{enumerate}[label=\alph*)]
\item Với $x\ge0$, $x\ne4$, ta có
\begin{align*}
A&=\frac{x}{x-4}+\frac{1}{\sqrt{x}-2}+\frac{1}{\sqrt{x}+2}\\
&=\frac{x+(\sqrt{x}+2)+(\sqrt{x}-2)}
{(\sqrt{x}+2)(\sqrt{x}-2)}\\
&=\frac{\sqrt{x}(\sqrt{x}+2)}
{(\sqrt{x}+2)(\sqrt{x}-2)}
=\frac{\sqrt{x}}{\sqrt{x}-2}.
\end{align*}
\item Với $x=25$,
\[
A=\frac5{5-2}=\frac53.
\]
\item
\[
\frac{\sqrt{x}}{\sqrt{x}-2}=-\frac13
\Longleftrightarrow 3\sqrt{x}=-\sqrt{x}+2
\Longleftrightarrow \sqrt{x}=\frac12
\Longleftrightarrow x=\frac14.
\]
Giá trị này thỏa mãn điều kiện.
\end{enumerate}

\textbf{Câu 2.}

Gọi số áo tổ II may trong một ngày là $x$ ($x\in\mathbb N^*$). Khi đó mỗi ngày tổ I may được $x+10$ chiếc áo. Theo đề bài,
\[
3(x+10)+5x=1310
\Longleftrightarrow 8x=1280
\Longleftrightarrow x=160.
\]
Vậy mỗi ngày tổ I may 170 chiếc áo và tổ II may 160 chiếc áo.

\textbf{Câu 3.}
\begin{enumerate}[label=\alph*)]
\item Khi $m=1$, phương trình trở thành
\[
x^2-4x+3=0
\Longleftrightarrow (x-1)(x-3)=0.
\]
Vậy $x=1$ hoặc $x=3$.

\item Phương trình có hai nghiệm phân biệt khi
\[
\Delta'=(m+1)^2-(m^2+2)=2m-1>0
\Longleftrightarrow m>\frac12.
\]
Theo định lý Viète,
\[
x_1+x_2=2(m+1),\qquad x_1x_2=m^2+2.
\]
Do đó
\begin{align*}
x_1^2+x_2^2=10
&\Longleftrightarrow (x_1+x_2)^2-2x_1x_2=10\\
&\Longleftrightarrow 4(m+1)^2-2(m^2+2)=10\\
&\Longleftrightarrow m^2+4m-5=0\\
&\Longleftrightarrow m=1\ \text{hoặc}\ m=-5.
\end{align*}
Đối chiếu điều kiện $m>\dfrac12$, ta được $\boxed{m=1}$.
\end{enumerate}

\textbf{Câu 5.}

Ta có
\[
\sqrt{x^2+x+\frac14}=\left|x+\frac12\right|
\]
và vế phải bằng
\[
\frac12(2x+1)(x^2+1)=\left(x+\frac12\right)(x^2+1).
\]
Để phương trình có nghiệm, vế phải phải không âm. Vì $x^2+1>0$, suy ra
$x\ge-\dfrac12$. Khi đó phương trình tương đương
\begin{align*}
\left(x-\frac12\right)\left(x+\frac12\right)+x+\frac12
&=\left(x+\frac12\right)(x^2+1).
\end{align*}
Suy ra
\begin{align*}
\left(x+\frac12\right)^2
&=\left(x+\frac12\right)(x^2+1)\\
\Longleftrightarrow
\left(x+\frac12\right)x^2&=0\\
\Longleftrightarrow x&=-\frac12\quad\text{hoặc}\quad x=0.
\end{align*}
Cả hai giá trị đều thỏa mãn phương trình. Vậy
\[
\boxed{S=\left\{-\frac12,0\right\}}.
\]
\end{document}
